%
% Halaman Abstrak
%
% @author  Andreas Febrian
% @version 2.2.0
% @edit by Ichlasul Affan
%

\chapter*{Abstrak}
\singlespacing

\noindent \begin{tabular}{l l p{10cm}}
	\ifx\blank\npmDua
		Nama&: & \penulisSatu \\
		Program Studi&: & \programSatu \\
	\else
		Nama Penulis 1 / Program Studi&: & \penulisSatu~/ \programSatu\\
		Nama Penulis 2 / Program Studi&: & \penulisDua~/ \programDua\\
	\fi
	\ifx\blank\npmTiga\else
		Nama Penulis 3 / Program Studi&: & \penulisTiga~/ \programTiga\\
	\fi
	Judul&: & \judul \\
	Pembimbing&: & \pembimbingSatu \\
	\ifx\blank\pembimbingDua
    \else
        \ &\ & \pembimbingDua \\
    \fi
    \ifx\blank\pembimbingTiga
    \else
    	\ &\ & \pembimbingTiga \\
    \fi
\end{tabular} \\

\vspace*{0.5cm}

Klaster Kubernetes yang dinamis mengakumulasi fragmentasi kapasitas CPU dan memori ketika beban kerja berubah. Walaupun kapasitas agregat masih mencukupi, bentuk sumber daya yang tidak kompatibel dapat menyebabkan kapasitas terdampar pada beberapa \textit{Node} dan menghambat penempatan beban kerja. Penelitian ini menyajikan arsitektur Kubernetes Descheduler multi-tingkat yang menggabungkan kebijakan \textit{ResourceDefragmentation} dengan penyaringan penggusuran sadar jaringan. Kebijakan sumber daya menilai \textit{Node} menggunakan utilisasi rata-rata, \textit{bin score} CPU-memori, \textit{Resource Imbalance Index} (RII), dan \textit{Free-Space Index} (FSI). Sebelum penggusuran, kebijakan memprediksi target penjadwal yang layak dan memilih korban berdasarkan proyeksi \textit{bin score} CPU-memori pada \textit{Node} tujuan. Penggusuran hanya dilakukan jika target tetap berada di bawah batas utilisasi dan setidaknya sama padatnya dengan sumber. Pada lima skenario berbasis \textit{request}, kebijakan mereklamasi empat \textit{Node} pada skenario utilisasi rendah, tiga \textit{Node} pada skenario terfragmentasi, serta menurunkan skor kapasitas terdampar dari \(1{,}338\) menjadi \(0{,}123\) pada skenario pengosongan heterogen dengan dua penggusuran. Seluruh eksperimen mencapai konvergensi tanpa meninggalkan \textit{Pod} dalam status \textit{Pending}. Lapisan sadar jaringan juga mencegah perpindahan yang memperburuk kedekatan komunikasi. Hasil ini menunjukkan bahwa keseimbangan target prediksi dapat menghasilkan defragmentasi yang efektif dengan tetap mempertahankan pengaman kelayakan dan lokalitas jaringan.

\vspace*{0.2cm}

\noindent Kata kunci: \\ Kubernetes, Descheduler, Resource Defragmentation, Latensi Jaringan, Microservices \\

\setstretch{1.4}
\newpage
